\begin{table}[htbp]
    \centering
    \caption{Planificación temporal del desarrollo}
    \label{tab:gantt_pipeline}
    \small
    \begin{tabular}{@{}p{5.5cm}cccccccc@{}}
    \toprule
    \textbf{Tarea / Iteración} &
    \textbf{S.1} & \textbf{S.2} & \textbf{S.3} & \textbf{S.4} &
    \textbf{S.5} & \textbf{S.6} & \textbf{S.7} & \textbf{S.8} \\
    \midrule
    1. Captura y envío de audio (móvil + \englishText{gateway}) &
    \cellcolor{blue!15} & \cellcolor{blue!15} & & & & & & \\
    2. Transcripción \englishText{ASR} por fragmentos &
    & \cellcolor{blue!15} & \cellcolor{blue!15} & & & & & \\
    3. \englishText{NER} dual y fusión de entidades &
    & & \cellcolor{blue!15} & \cellcolor{blue!15} & & & & \\
    4. Enriquecimiento \englishText{SNOMED CT} y \englishText{CIE-10-ES} &
    & & & \cellcolor{blue!15} & \cellcolor{blue!15} & & & \\
    5. Reporte de control e informe \englishText{LLM} &
    & & & & \cellcolor{blue!15} & \cellcolor{blue!15} & & \\
    6. \englishText{Bundle FHIR} R4B &
    & & & & & \cellcolor{blue!15} & \cellcolor{blue!15} & \\
    7. Validación \englishText{end-to-end} y redacción del TFM &
    & & & & & & \cellcolor{blue!15} & \cellcolor{blue!15} \\
    \bottomrule
    \end{tabular}
    \end{table}
    Las columnas S.1--S.8 representan sprints de desarrollo consecutivos del periodo de desarrollo del trabajo fin de
    máster. Cada sprint tiene una duración de 2 semanas. Las celdas sombreadas indican el intervalo de ejecución de cada iteración.